\documentclass[12pt]{report}
\usepackage{amsmath}
\usepackage{amssymb}
\usepackage{amsthm}
\usepackage{enumitem}
\usepackage{multicol}

\newenvironment{case}[2]{\par\textbf{Case #1: #2.}}{}

\begin{document}
\renewcommand\qedsymbol{OE$\Delta$}
\newcommand{\abs}[1]{\left| #1 \right|}
\newcommand{\andtxt}{\text{ and }}
\newcommand{\Int}{\text{Int}}
\newcommand{\Ext}{\text{Ext}}
\newcommand{\mathif}{\text{if }}
\newcommand{\seq}[2][n]{\ensuremath{(#2 _{#1})}}

\tableofcontents
\pagebreak

\setcounter{chapter}{5}
\chapter{The Continuity Struggle}
\pagebreak

\setcounter{section}{11}
\section{Problem 6.12}
    Let $f: \mathbb{R} \to \mathbb{R}$. 
    Prove that $\lim_{x \to 0} f\left( \frac{1}{x} \right) = \lim_{x \to 1} f(x)$. 
    That is, prove that if one converges to some $L$, then the  other also converges to $L$; if one diverges to $\infty$ or to $-\infty$ or whose limit does not exist, then the other does the same. 

    \subsection{Solution}
    \begin{proof}[Proof by Contradiction]
        Suppose for contradiction that the two limits are not equal.
        \begin{gather}
            \lim_{x \to 0^+} f\left( \frac{1}{x} \right) \neq \lim_{x \to \infty} f(x)
        \end{gather}

        Suppose there exists a function $g: \mathbb{R} \to \mathbb{R}$ defined aas taking the reciprocal of its initial value.
        \begin{equation}
            g(x) := \frac{1}{x}
        \end{equation}

        We can prove that $g(x)$ converges to $0$ from the positive direction as $x$ approaches $\infty$.
        Suppose Let $\varepsilon > 0$ and let $\delta = \frac{1}{\epsilon} > 0$.
        Then, for any $x$ that satisfies $x > \delta$ and $x > 0$, we can take its reciprocal.
        \begin{gather}
            0 < \frac{1}{x} < \frac{1}{\delta}
        \end{gather}
        
        Recall that, by the sequence squeeze theorem and the Archimedian Principle, for every sequence $x_n$ that converges to $\infty$, $\frac{1}{x_n}$ converges to $0$, so we can replace $\frac{1}{\infty}$ with $0$.
        \begin{align}
            \abs{g(x) - g(\infty)}   &=  \abs{\frac{1}{x} - \frac{1}{\infty}}\\
                &=  \abs{\frac{1}{x} - 0}
                =   \abs{\frac{1}{x}}
        \end{align}
        
        This has an attached inequality from earlier.
        \begin{gather}
            \abs{g(x) - g(\infty)}  <  \abs{\frac{1}{\delta}}
                =   \abs{\frac{1}{\frac{1}{\varepsilon}}}
                =   \abs{\varepsilon}
                =   \varepsilon\\
            \abs{g(x) - g(\infty)} < \varepsilon
        \end{gather}

        Therefore, $g(x)$ converges to $0$ from the positive direction as $x$ approaches $\infty$.

        Suppose we made a sequence $a_n$ that converges to $0$ and another sequence $b_n$ that converges to $0$ from the positive direction as $n$ approaches $\infty$.
        $g(b_n)$ would in turn diverge to $\infty$.
        We can rewrite our initial supposition.
        \begin{gather}
            \lim_{x \to 0^+} f\left( \frac{1}{x} \right) \neq \lim_{x \to \infty} f(x)
        \end{gather}

        We can rewrite the latter term by substituting in $g(x)$.
        \begin{equation}
            \lim_{x \to 0^+} f\left( \frac{1}{x} \right) = \lim_{n \to \infty} \frac{1}{g(b_n)}
                =   \lim_{n \to \infty} \frac{1}{\frac{1}{b_n}}
                =   \lim_{n \to \infty} b_n
        \end{equation}

        We can rewrite the former term by substitutring in $a_n$.
        \begin{equation}
            \lim_{x \to \infty} f(x) = \lim_{n \to \infty} a_n
        \end{equation}

        Apply these to the inequality.
        \begin{equation}
            \lim_{n \to \infty} b_n \neq \lim_{n \to \infty} a_n
        \end{equation}

        This contradicts the idea that functional limits are unique.
        Ergo, our initial assumption is false and $\lim_{x \to 0} f\left( \frac{1}{x} \right) = \lim_{x \to 1} f(x)$. 
    \end{proof}

\end{document}